\section{Discussion}
\label{sec:conclusion}

In this work, we argue that verification constitutes an underexplored axis of scaling. To realize this, we propose LLM-as-a-Verifier, a general-purpose framework that delivers fine-grained feedback for agentic tasks. Unlike standard LM judges that output a single discrete score, our approach computes a continuous reward by taking the expectation over the distribution of scoring-token logits and enables verification scaling across multiple dimensions, including (1) score granularity, (2) repeated evaluation, and (3) criteria decomposition. When used as a trajectory reward model for test-time scaling, it achieves state-of-the-art performance on Terminal-Bench V2, SWE-Bench Verified, RoboRewardBench, and MedAgentBench. Beyond ranking, the fine-grained verifier signal can be used as a progress estimator, opening a path toward safer real-world deployment of autonomous agents. Finally, we demonstrate that LLM-as-a-Verifier can be used as a dense reward signal for RL, improving the sample efficiency of SAC and GRPO on robotics and mathematical reasoning benchmarks.

\clearpage

\section{Related Work}
\label{app:related-work}

\paragraph{Test-Time Scaling.}
Test-time scaling improves model performance by spending additional inference compute on deliberation, search, or candidate generation. One line of work improves a single response by eliciting intermediate reasoning with chain-of-thought~\citep{wei_chain--thought_nodate,kojima2022large}, decomposing problems into simpler subproblems~\citep{zhou2022least}, or marginalizing over sampled reasoning paths~\citep{wang_self-consistency_2023}. Another line searches over intermediate thoughts~\citep{yao2023tree,besta2023graph}, actions~\citep{yao_react_2023,zhou2023language}, or latent world states~\citep{hao2023rap} based on test-time feedback~\citep{weng2022selfverification,madaan_self-refine_nodate,shinn_reflexion_nodate,gou2023critic,agrawal_gepa_2025,novikov_alphaevolve_2025}. Repeated sampling and best-of-$N$ selection further scale candidate pools for code generation~\citep{li2022alphacode}, general reasoning~\citep{brown2024largelanguagemonkeysscaling,snell_scaling_2024}, parallel self-verification~\citep{singh_v_1_2026}, and inference-aware training~\citep{chow2024inferenceaware}. These approaches can expose substantial oracle headroom, but realizing this headroom requires a reliable selector. Our LLM-as-a-Verifier instead shows that verifier quality can be improved by scaling score granularity, repeated evaluation, and criteria decomposition, and that better verification directly improves best-of-$N$ selection for long-horizon decision making.

\paragraph{LLM-as-a-Judge.}
LLM-as-a-judge methods provide a scalable alternative to human evaluation by prompting large models to score or compare generated outputs. Probability- and form-filling-based evaluators extract richer scoring signals from LLMs~\citep{fu2023gptscore,liu2023geval}, while benchmark-style evaluators use LLM preferences to evaluate instruction-following systems~\citep{zheng2023judging,dubois2024alpacaeval,li2024arenahard}. A complementary line makes evaluation more fine-grained through skill decompositions~\citep{ye2023flask}, customized rubrics and specialized open judges~\citep{kim2024prometheus,kim2024prometheus2,li2024autoj,hu2024themis}, and hierarchical criteria~\citep{liu2024hdeval}. Other work trains scalable judge models~\citep{wang2024pandalm,zhu2025judgelm} or combines multiple judges through debate and weak-verifier ensembling~\citep{chan2023chateval,saad2025shrinking}. Recent studies have also analyzed judge reliability, including fairness and position biases~\citep{wang2024large,zeng2024llmbar}, cognitive and self-enhancement biases~\citep{koo2023cobbler,liu2023narcissistic}, general judge benchmarks~\citep{tan2025judgebench,huang2025empirical}, as well as domain-specific judge evaluation~\citep{jiang2025codejudgebench}. Multimodal judge models extend this paradigm to image and vision-language evaluation~\citep{chen2024mllmjudge,lee2024prometheusvision,xiong2024llavacritic}. Our work builds on these works but differs in setting, objective, and scaling characterization. Rather than evaluating isolated natural-language responses, LLM-as-a-Verifier verifies long-horizon agent trajectories involving tool use, code execution, robotics, and medical decision-making. Moreover, we systematically study how verification quality scales with score granularity, repeated evaluations, and criteria decomposition. 

\paragraph{Verifiable Reward.}
Reward models convert candidate solutions, actions, or trajectories into scalar feedback for selection, monitoring, or policy optimization. In language reasoning, learned verifiers have been trained as outcome reward models for final-answer selection~\citep{cobbe_training_2021}, process reward models for step-level supervision~\citep{uesato2022solving,lightman_lets_2023}, and generative verifiers that cast reward modeling as next-token prediction~\citep{zhang2025generativeverifiersrewardmodeling}. 
In robotics, reward signals have been derived from value-implicit visual representations~\citep{ma2022vip}, language-image reward representations~\citep{ma2023liv}, pretrained vision-language models~\citep{sontakke2023roboclip,rocamonde2023vlmrm,ICLR2025_54854cf1}, VLM feedback or preferences~\citep{wang2024rlvlmf}, LLM-generated reward code~\citep{yu2023languagetorewards,xie2023text2reward,ma2023eureka}, token-probability progress~\citep{chen2026topreward}, and trained general-purpose robotic reward models based on large-scale trajectory or preference data~\citep{zhang2025rewind,chen2025sarm,lee2026roboreward,liang2026robometer}.
Recent work has also developed action-level verifiers for guided sampling~\citep{nakamoto2024steering,liu2024bidirectional,kwok2025robomonkey}, runtime monitoring~\citep{agia2024unpacking}, and multimodal alignment~\citep{kwok2026scalingverificationeffectivescaling}.
A complementary line makes verification more reliable by factorizing holistic judgments into smaller checks~\citep{min2023factscore,fabbri2021qafacteval,manakul2023selfcheckgpt,liu2024hdeval,liu2026worldactionverifierselfimproving,tseng2026sc3}. Orthogonal to these methods, our work studies how verifier quality scales with score granularity, repeated evaluation, and criteria decomposition across multiple domains in a general-purpose framework.


% \paragraph{Test-Time Scaling.}
% Test-time scaling has emerged as an effective paradigm for improving LLM performance without additional training, including longer reasoning traces ~\citep{wei_chain--thought_nodate}, self-refinement~\citep{shinn_reflexion_nodate}, search~\citep{yao2023tree,zhou2023language}, and best-of-$N$ sampling~\citep{brown2024largelanguagemonkeysscaling,liu2024bidirectional}. These methods~\citep{agrawal_gepa_2025, novikov_alphaevolve_2025, brown2024largelanguagemonkeysscaling, snell_scaling_2024, singh_v_1_2026} improve generation by allocating more inference compute to either a single trajectory or a larger pool of candidate solutions. However, they often assume access to a reliable evaluator, such as a deterministic checker, majority vote, or coarse LM judge. In contrast, our work studies verification itself as the object of scaling. We show that candidate pools often contain substantial oracle headroom, and that downstream performance is bottlenecked by verification rather than by generation alone.

% \paragraph{LLM-as-a-Judge.}
% LLM-as-a-judge methods provide a scalable alternative to human evaluation by prompting LLMs to assess generated outputs~\citep{liu_g-eval_2023,li2024autoj}. Prior work such as MT-Bench and Chatbot Arena~\citep{zheng_judging_nodate} established the effectiveness of LLM judges for open-ended response evaluation, while G-Eval~\citep{liu_g-eval_2023} introduced chain-of-thought reasoning and probability-weighted scoring to improve alignment with human preferences. Subsequent work has studied fine-grained, rubric-conditioned, and trained judge models~\citep{ye2023flask,li2024autoj}, as well as multi-judge or ensemble-based evaluation~\citep{chan2023chateval,saad2025shrinking}. Weaver~\citep{saad2025shrinking} further shows that ensembling weak reward models can reduce the generation-verification gap. Our work builds on this line of work but differs in setting, objective, and scaling characterization. Rather than evaluating isolated natural-language responses, LLM-as-a-Verifier verifies long-horizon agent trajectories involving tool use, code execution, robotics, and medical decision-making. Moreover, instead of collapsing the model's scoring distribution into a single discrete rating, we use a probabilistic approach and systematically study how verification quality scales with score granularity, repeated evaluations, and criteria decomposition. 

% \paragraph{Reward Models.}
% Learned reward models and verifiers~\citep{cobbe_training_2021, zhang2025generativeverifiersrewardmodeling} have been widely used to improve reasoning and decision-making. Outcome reward models score final solutions, while process reward models provide supervision over reasoning steps~\citep{uesato2022solving,lightman_lets_2023}. In robotics, learned reward models and action verifiers~\citep{kwok_robomonkey_2025, kwok2026scalingverificationeffectivescaling, liu2026worldactionverifierselfimproving} have been used to score candidate actions or trajectories. However, these approaches typically require domain-specific training data and may not generalize across tasks with different modalities, horizons, or success criteria. LLM-as-a-Verifier instead uses a frozen language or vision-language model as a general-purpose trajectory reward model. The same approach can be applied across different domains without additional training.